\section{Cahier des charges}

\subsection{Contexte et besoin}

Le besoin identifié est de permettre à un utilisateur d'identifier une musique
diffusée autour de lui sans utiliser directement son téléphone. L'objet doit
pouvoir être posé dans un environnement sonore, lancé par un bouton physique,
puis afficher le résultat de manière claire.

Le prototype s'adresse principalement à :

\begin{itemize}
    \item un utilisateur curieux souhaitant identifier rapidement un morceau ;
    \item un maker voulant reproduire ou modifier le système ;
    \item un jury ou un public de démonstration ;
    \item un développeur souhaitant expérimenter la reconnaissance audio sur
    Raspberry Pi.
\end{itemize}

\subsection{Objectif principal}

Réaliser un prototype fonctionnel capable de capturer un extrait audio,
d'envoyer cet extrait à un service de reconnaissance musicale, puis d'afficher
le résultat sur un écran e-paper.

\subsection{Objectifs secondaires}

\begin{itemize}
    \item fournir une interface web locale pour la configuration Wi-Fi ;
    \item afficher un QR code de configuration en mode setup ;
    \item gérer des boutons physiques pour piloter le prototype ;
    \item documenter l'architecture logicielle et matérielle ;
    \item préparer une base de code modulaire et évolutive ;
    \item garder une solution compatible avec les ressources d'un Raspberry Pi.
\end{itemize}

\subsection{Périmètre}

\begin{longtable}{|p{0.45\linewidth}|p{0.45\linewidth}|}
\hline
\textbf{Inclus dans le prototype} & \textbf{Hors périmètre initial} \\
\hline
Capture audio courte via ALSA et \texttt{arecord}. &
Production industrielle du boîtier. \\
\hline
Reconnaissance musicale via API AudD. &
Reconnaissance musicale 100 \% locale. \\
\hline
Affichage e-paper Waveshare 4.2 pouces V2. &
Application mobile native. \\
\hline
Interface web Flask pour configuration Wi-Fi. &
Historique musical avancé multi-utilisateur. \\
\hline
Boutons physiques sur GPIO. &
Optimisation poussée de l'autonomie. \\
\hline
Structure de documentation et code Python. &
Certification ou conformité industrielle. \\
\hline
\end{longtable}

\subsection{Contraintes}

\subsubsection{Contraintes matérielles}

\begin{itemize}
    \item utilisation d'un Raspberry Pi Zero 2 W ;
    \item intégration d'un microphone MEMS ICS-43434 en I2S ;
    \item intégration d'un écran Pico-ePaper-4.2 de 400 par 300 pixels ;
    \item utilisation de boutons physiques ;
    \item alimentation par batterie LiPo 2000 mAh, 7,4 Wh, avec UPS HAT (C) ;
    \item boîtier ou support inspiré d'une cassette futuriste ;
    \item coût compatible avec un prototype étudiant.
\end{itemize}

\subsubsection{Contraintes logicielles}

\begin{itemize}
    \item langage principal : Python ;
    \item capture audio par commande système \texttt{arecord} ;
    \item serveur web local avec Flask ;
    \item gestion Wi-Fi via \texttt{nmcli} ;
    \item affichage via la bibliothèque e-paper Waveshare ;
    \item code organisé en modules indépendants.
\end{itemize}

\subsubsection{Contraintes de délai}

Le projet a été conduit du 11 mai 2026 au 11 juin 2026. La durée disponible a
imposé une démarche de prototypage incrémental : valider d'abord la chaîne
fonctionnelle principale, puis améliorer l'ergonomie et la documentation.

\subsection{Fonctions attendues}

\begin{longtable}{|p{1.5cm}|p{8.8cm}|p{2.2cm}|}
\hline
\textbf{Code} & \textbf{Fonction attendue} & \textbf{Priorité} \\
\hline
F1 & Capturer un extrait audio exploitable. & Haute \\
\hline
F2 & Lancer une reconnaissance musicale à partir de l'extrait capturé. & Haute \\
\hline
F3 & Récupérer les métadonnées disponibles : titre, artiste, album, date, lien. & Haute \\
\hline
F4 & Afficher un état clair sur l'écran e-paper. & Haute \\
\hline
F5 & Afficher le résultat reconnu sur l'écran e-paper. & Haute \\
\hline
F6 & Afficher un QR code pour accéder à la configuration Wi-Fi. & Moyenne \\
\hline
F7 & Fournir une interface web locale de configuration. & Haute \\
\hline
F8 & Piloter les actions principales par boutons physiques. & Haute \\
\hline
F9 & Gérer les erreurs sans bloquer l'utilisateur. & Haute \\
\hline
F10 & Documenter le montage, les choix et les tests. & Haute \\
\hline
\end{longtable}

\subsection{Critères de validation}

\begin{longtable}{|p{4cm}|p{9cm}|}
\hline
\textbf{Élément} & \textbf{Critère de validation} \\
\hline
Capture audio & Un fichier WAV temporaire est généré après appui sur le bouton d'action. \\
\hline
Reconnaissance & Un morceau audible et connu retourne au minimum un titre et un artiste lorsque l'API répond. \\
\hline
Affichage & L'écran e-paper affiche les états \textit{idle}, \textit{listening}, \textit{thinking} et résultat. \\
\hline
Wi-Fi & L'interface web permet de scanner les réseaux et d'envoyer un SSID/mot de passe. \\
\hline
Boutons & Le bouton de mode change entre setup et running ; le bouton d'action lance une écoute manuelle. Chaque nouvelle écoute nécessite un nouvel appui. \\
\hline
Documentation & Le dossier décrit le besoin, les fonctions, l'architecture, le planning et les tests. \\
\hline
\end{longtable}
